%------------------------------------------------------------------------------%
\addtocounter{exemplo}{1}
\begin{frame}{Exemplo \theexemplo}
  Encontre e classifique os pontos críticos da função
  \[
    f(x, y) = 3y^2 -2y^3 -3x^2 +6xy
  \]
  Avalie os valores extremos locais de $f$
\end{frame}

%------------------------------------------------------------------------------%
\begin{frame}{Exemplo \theexemplo}
  Derivadas parciais
  \pause
  \vspace{\baselineskip}
  \[
    \D{f}{x}
    \pause
    = \D{}{x}\left(3y^2 -2y^3 -3x^2 +6xy\right)
    \pause
    = -6x +6y
    \pause
    = 6y -6x
  \]
  \pause
  \vspace{\baselineskip}
  \[
    \D{f}{y}
    \pause
    = \D{}{y}\left(3y^2 -2y^3 -3x^2 +6xy\right)
    \pause
    = 6y -6y^2 +6x
  \]
\end{frame}

%------------------------------------------------------------------------------%
\begin{frame}{Exemplo \theexemplo}
  Pontos críticos

  \pause
  \vspace{\baselineskip}
  As derivadas existem em todos os pontos do plano
\end{frame}

%------------------------------------------------------------------------------%
\begin{frame}{Exemplo \theexemplo}
\begin{columns}
\column{0.5\linewidth}
  \begin{align*}
    \onslide<+->{f_x(x, y) & = 6y -6x = 0 \\}
    \onslide<+->{f_y(x, y) & = 6y -6y^2 +6x = 0}
  \end{align*}
  \begin{align*}
    \onslide<+->{6y -6x & = 0 \\}
    \onslide<+->{x      & = y}
  \end{align*}
  \column{0.5\linewidth}
  \begin{align*}
    \onslide<+->{6y -6y^2 +6x & = 0 \\}
    \onslide<+->{6y^2 -6y -6x & = 0 \\}
    \onslide<+->{6y^2 -6y -6y & = 0 \\}
    \onslide<+->{6y^2 - 12y   & = 0 \\}
    \onslide<+->{6y(y-2)      & = 0}
  \end{align*}

  \vspace{\baselineskip}
  \onslide<+->{Portanto \(y=0\) ou \(y=2\)}

  \vspace{\baselineskip}
  \onslide<+->{Os pontos críticos são \((0, 0)\) e \((2, 2)\)}
\end{columns}
\end{frame}

%------------------------------------------------------------------------------%
\begin{frame}{Exemplo \theexemplo}
\begin{columns}[t]
\column{0.5\linewidth}
  \onslide<+->{Calculando as derivadas segundas}
  \[
    \onslide<+->{\DS{f}{x}}
    \onslide<+->{= \D{}{x}\left(6y -6x\right)}
    \onslide<+->{= -6}
  \]

  ~
  \[
    \onslide<+->{\DS{f}{y}}
    \onslide<+->{= \D{}{y}\left(6y -6y^2 +6x\right)}
    \onslide<+->{= 6 -12y}
  \]

  ~
  \[
    \onslide<+->{\DM{f}{x}{y}}
    \onslide<+->{= \DM{f}{y}{x}}
    \onslide<+->{= \D{}{y}\left(6y -6x\right)}
    \onslide<+->{= 6}
  \]
\column{0.5\linewidth}
  \onslide<+->{Discriminante}
  \begin{align*}
    \onslide<+->{D(x, y)}
    \onslide<+->{& = f_{xx}(x, y) f_{yy}(x, y) - f^2_{xy}(x, y) \\}
    \onslide<+->{& = (-6)(6-12y) - 6^2 \\}
    \onslide<+->{& = -36 +72y - 36 \\}
    \onslide<+->{& = 72(y-1)}
  \end{align*}
\end{columns}
\end{frame}

%------------------------------------------------------------------------------%
\begin{frame}{Exemplo \theexemplo}
  Analisando o ponto \((0, 0)\)

  \pause
  \vspace{\baselineskip}
  Como
  \[
    \pause
    D(0, 0)
    \pause
    = 72(0-1)
    \pause
    = -72
    \pause
    \;<\; 0
  \]
  \pause
  esse ponto é um ponto de sela
\end{frame}

%------------------------------------------------------------------------------%
\begin{frame}{Exemplo \theexemplo}
  \onslide<+->{Analisando o ponto \((2, 2)\)}

  \onslide<+->{Como}
  \[
    \onslide<+->{D(2, 2)}
    \onslide<+->{= 72(2-1)}
    \onslide<+->{= 72}
    \onslide<+->{> 0}
    \qquad
    \onslide<+->{f_{xx}}
    \onslide<+->{= -6}
    \onslide<+->{< 0}
  \]
  \onslide<+->{o ponto \((2, 2)\) é um máximo local com valor}
  \begin{align*}
    \onslide<+->{f(2, 2)}
    \onslide<+->{& = \left(3y^2 -2y^3 -3x^2 +6xy\right)\evalat{(2, 2)}{} \\}
    \onslide<+->{& = 3\times 2^2 -2 \times 2^3 -3\times 2^2 + 6\times 2 \times 2\\}
    \onslide<+->{& = 12 -16 -12 + 24 \\}
    \onslide<+->{& = 8}
  \end{align*}
\end{frame}

%------------------------------------------------------------------------------%
